\documentclass[report.tex]{subfiles}

\begin{document}

\section{Verified Compilation Background}

\subsection{Small-Step Operational Semantics}

Giving a formal semantics to a programming language
provides a rigorous meaning to its programs.
Small-step operational semantics models computation
as a sequence of elementary transitions between program configurations.
Rather than evaluating compound statements in a single step,
it describes how internal configurations evolve under primitive operations.
This style directly captures non-termination as infinite transition sequences
and permits the inspection of intermediate configurations.

Formally,
a small-step operational semantics is defined by a transition system
over a set of configurations $C$,
equipped with a binary reduction relation
$\twoheadrightarrow \;\subseteq C \times C$.
In our framework,
a configuration $c$ is represented as a pair $c = (s, m) \in S \times M$,
where $s \in S$ denotes an internal program state
and $m \in M$ denotes a memory state.
The structure of the internal state $s$ depends on the language under study:
for instance,
$s$ may be a term in a $\lambda$-calculus with substitution semantics,
or,
as in the \muRTL{} language introduced later in this report,
a program,
a call stack,
a currently active function,
a program counter,
and an environment.

The transition relation $\step{(s, m)}{(s', m')}$
indicates that the configuration $(s, m)$ reduces in a single step
to $(s', m')$.
The global program code and declarations are implicitly fixed
within the transition relation.
We denote by $\steps{}{}$
the reflexive-transitive closure of $\step{}{}$.
The transition system need not be deterministic:
a configuration may admit several outgoing transitions.

To characterize execution endpoints,
the semantics is parameterized by a set of observable values $\val$
and a partial function \isfinalf.
A configuration $c$ is \emph{final}
when $\isfinal{c} = \Some{v, m}$ for some return value $v \in \val$
and memory $m \in \mem$.
For such final configurations,
no further reduction steps are possible:
if $\isfinal{c} = \Some{v, m}$,
there exists no configuration $c'$
such that $\step{c}{c'}$.

We say that a configuration $c$ can make progress,
written $\canp{c}$,
if there exists a configuration $c'$
such that $\step{c}{c'}$.
Conversely,
we follow the convention that a configuration $c$ is \emph{stuck},
written $\stuck{c}$,
if it is not a final configuration
and cannot perform any transition.
Stuck configurations represent execution failures
where the operational semantics is undefined.
In an imperative setting,
a configuration becomes stuck when dereferencing an unallocated pointer,
dividing by zero,
or branching to an invalid control-flow node.

\subsection{Program Behaviors}

Equipped with a small-step operational semantics,
we can formally classify the observable behaviors of a program execution.
Because the transition system may be non-deterministic,
an execution path starting from a configuration $c$
can yield multiple outcomes.
The syntax of observable behaviors $b$ is defined by the following grammar:
\begin{equation*}
  b \in \beh \Coloneqq \Term{v, m} \mid \Div \mid \Undef
\end{equation*}
where $v \in \val$ and $m \in \mem$.
We write $b \in c$
to say that the configuration $c$ exhibits the behavior $b$.
We say that:
\begin{itemize}
  \item $c$ \emph{terminates} with value $v$ and memory $m$
        if there is a final configuration $c'$
        such that $\steps{c}{c'}$ and $\isfinal{c'} = \Some{v, m}$.

  \item $c$ \emph{diverges} if there exists an infinite reduction sequence
        starting from it.
        We do not restrict divergence to strongly diverging configurations,
        where every execution path is infinite.

  \item $c$ exhibits \emph{undefined behavior}
        if there is a configuration $c'$
        such that $\steps{c}{c'}$ and $\stuck{c'}$.
\end{itemize}

The notion of behavior is extended to a complete program ($b \in P$)
by considering the behaviors of its initial state.
Every configuration admits at least one observable outcome:
\begin{lemma}
For any configuration $c$,
there exists a behavior $b$ such that $b \in c$.
\end{lemma}

\subsection{Program Refinement}

The objective of a verified compiler is to guarantee semantic preservation:
the compiled target executable must faithfully preserve
the behaviors of the source program.
In verified compilation,
semantic preservation is formalized as a \emph{behavioral refinement}
between the compiled target program $T$ and the source program $S$.

A direct formulation is the \emph{rigid behavioral refinement},
requiring every behavior of the target program
to be an admitted behavior of the source program:
\begin{equation}
  \label{eq:rigid-beh-ref}
  T \sqsubseteq S \triangleq \forall b, b \in T \implies b \in S
\end{equation}
Rigid refinement ensures that compilation never introduces spurious behaviors,
while permitting the compiler to resolve non-deterministic choices
present in the source.
However,
the rigid behavioral refinement of \cref{eq:rigid-beh-ref}
is overly restrictive for optimizing compilers.
If a source program only exhibits undefined behaviors,
an optimizing compiler may eliminate an erroneous code path
or transform it into a benign execution for the target.

To account for these optimization principles,
compiler correctness relies on a \emph{flexible behavioral refinement}.
We first equip behaviors with a preorder $\behord$,
parameterized by a relation on values and memory $\Phi$:
\begin{gather*}
  \begin{prooftree}
    \hypo{\Phi(v_t, v_s, m_t, m_s)}
    \infer1[\textsc{Beh-Term}]{
      \Term{v_t, m_t} \behord \Term{v_s, m_s}
    }
  \end{prooftree}
  \qquad\qquad
  \begin{prooftree}
    \hypo{}
    \infer1[\textsc{Beh-Div}]{
      \Div \behord \Div
    }
  \end{prooftree}
  \\[1em]
  \begin{prooftree}
    \hypo{}
    \infer1[\textsc{Beh-Undef}]{
      b_t \behord \Undef
    }
  \end{prooftree}
\end{gather*}

These rules can be read as follows:
\begin{itemize}
  \item Producing a value $v_t$ and a memory $m_t$ in the compiled program
        is justified by the possibility of producing a value
        $v_s$ and a memory $m_s$ in the source program such that
        $\Phi(v_t, v_s, m_t, m_s)$.

  \item Divergence in the compiled program
        is justified by the possibility of divergence in the source program.

  \item Any behavior $b_t$ in the compiled program
        is justified by the possibility of undefined behavior in the source
        program.
\end{itemize}

We then define the \emph{flexible behavioral refinement}
between a target program $T$ and a source program $S$:
\begin{equation}
  \label{eq:flex-beh-ref}
  T \slantsqsubseteq S
  \triangleq
  \forall b_t, b_t \in T \implies
  \exists b_s,
  b_s \in S \land b_t \behord b_s
\end{equation}
This relation is similarly defined
between configurations $t \slantsqsubseteq s$.
Whenever $T \slantsqsubseteq S$ holds,
every observable behavior produced by the compiled program $T$
is validated by a matching behavior in the source program $S$
with respect to the preorder $\behord$.

\subsection{Simulation Relations}

While a refinement relation
is formulated globally over whole-program behaviors,
proving it directly over full execution traces is impractical.
Compilers perform transformations whose correctness is best verified
step by step.
To bridge this gap,
verified compilation relies on \emph{simulation relations},
which maintain an invariant relating intermediate configurations
of the source and target executions.
Simulation relations are established in two flavors:

\begin{description}
  \item[Backward simulation:]
        each execution step taken by the target program
        must be matched by corresponding steps in the source program.
        Because it maps every target execution path back to an admissible
        source path,
        a backward simulation directly establishes behavioral refinement
        $t \sqsubseteq s$.

  \item[Forward simulation:]
        each execution step taken by the source program
        is matched by steps in the target program.
        Forward simulations are more convenient to establish
        as compiler passes transform source constructs into target code,
        making inductive invariant proofs follow the structure of source
        reductions directly.
\end{description}

\subsection{Coinductive Predicates}

Proving a simulation relation requires relating execution steps
of the target and source programs.
When programs terminate,
their execution consists of a finite sequence of reductions ending in a
final configuration.
Such finite behaviors are formalized using \emph{inductive} predicates.
However,
programs that are designed to run indefinitely produce
infinite reduction sequences
that never encounter a terminal case.
Establishing a simulation inductively for diverging executions is impossible,
as it would require an infinite derivation tree,
which inductive definitions forbid.
To reason about infinite execution traces,
we rely on \emph{coinductive} predicates.

The presentation adopted below,
which interprets coinduction in terms of infinite derivation trees,
is an intuitive presentation adapted to our setting.
For foundational treatments of coinduction,
we refer the reader to
\textcites{milner-89-communication}{sangiorgi-12-bisimulation}%
{gimenez-98-tutorial}{pous-16-cawu}.
Throughout this report,
we specify inductive and coinductive relations through inference rules:
\begin{itemize}
  \item \emph{Inductive rules} are denoted with a single horizontal bar:
        \begin{equation*}
          \begin{prooftree}
            \hypo{H_1}
            \hypo{\dots}
            \hypo{H_n}
            \infer3{\mathcal{I}}
          \end{prooftree}
        \end{equation*}
        Here,
        $H_i$ are the premises and $\mathcal{I}$ is the inductive conclusion.
        A judgment holds inductively if and only if it can be established
        through a \emph{finite derivation tree},
        starting from base cases.

  \item \emph{Coinductive rules} are denoted with a double horizontal bar:
        \begin{equation*}
          \begin{prooftree}
            \hypo{H_1}
            \hypo{\dots}
            \hypo{H_n}
            \infer[coind]3{\mathcal{C}}
          \end{prooftree}
        \end{equation*}
        Here,
        $H_i$ are the premises and $\mathcal{C}$ is the coinductive conclusion.
        They generalize inductive rules to judgments admitting
        \emph{potentially infinite derivation trees}.
\end{itemize}
These inference rules can be read in a goal-directed manner:
\begin{equation*}
  \begin{prooftree}
    \hypo{P}
    \hypo{\text{Goal}'}
    \infer2{\text{Goal}}
  \end{prooftree}
\end{equation*}
To establish $\text{Goal}$,
one validates the side conditions $P$
and updates the proof goal to $\text{Goal}'$.
When $\text{Goal}'$ is an implication $H \implies \text{Goal}''$,
the hypothesis $H$ is assumed in the context to establish $\text{Goal}''$.
To combine inductive and coinductive reasoning within the same definition,
a relation $R$ is defined as the greatest fixed point%
\footnote{The relationship between greatest fixed points and coinduction
  follows from the Knaster--Tarski theorem;
  see \textcites{pous-16-cawu}{schafer-17-tower} for comprehensive treatments.}
of an inductive operator.
Inductive rules with a single bar require a finite proof tree
for occurrences of $R$ in the premises,
while coinductive rules with a double bar permit infinite derivations.

For example,
finite lists are defined inductively,
requiring every list to terminate at $\mathsf{nil}$.
In contrast,
infinite streams are defined coinductively,
allowing endless constructor applications without base cases:
\begin{equation*}
  \begin{prooftree}
    \hypo{}
    \infer1{\mathsf{nil} \in \List{\mathbb{N}}}
  \end{prooftree}
  \qquad
  \begin{prooftree}
    \hypo{n \in \mathbb{N}}
    \hypo{l \in \List{\mathbb{N}}}
    \infer2{\mathsf{cons}(n, l) \in \List{\mathbb{N}}}
  \end{prooftree}
  \qquad\qquad
  \begin{prooftree}
    \hypo{n \in \mathbb{N}}
    \hypo{s \in \Stream{\mathbb{N}}}
    \infer[coind]2{\mathsf{scons}(n, s) \in \Stream{\mathbb{N}}}
  \end{prooftree}
\end{equation*}

These distinct derivation structures
translate into distinct proof principles:
\begin{itemize}
  \item One reasons by \emph{induction}
        when the \emph{hypothesis} is an inductive judgment.
        The proof proceeds by case analysis
        on the constructors of the finite derivation.
        The induction hypothesis may only be applied to
        \emph{strictly smaller subterms},
        guaranteeing finite descent toward a base case.

  \item One reasons by \emph{coinduction}
        when the \emph{goal} to be established is a coinductive judgment.
        The proof constructs a potentially infinite derivation tree.
        The coinductive hypothesis refers to the goal itself,
        but it may only be invoked \emph{underneath a constructor}
        ($\mathsf{scons}$ for streams).
        This condition,
        known as \emph{productivity},
        prohibits using the coinductive hypothesis to build circular proof
        trees.
\end{itemize}

A detailed example illustrating coinductive reasoning is provided in
\cref{app:stream-bisimilarity}.
The Rocq theorem prover~\parencite{rocq}
supports coinductive predicates natively.
In our formalization,
we use the \texttt{coinduction} library developed by
\textcite{rocq-coinduction}
to facilitate reasoning with coinductive simulation invariants.

\subsection{Simulation Schemes \& Stuttering}

Equipped with coinductive predicates to handle potentially infinite
executions,
we can define a simulation relation to establish
program refinements: $T \slantsqsubseteq S$.
In general,
the source and target languages may differ:
target configurations $t$ evolve under $\step[t]{}{}$
with finality test $\isfinalf[t]$,
while source configurations $s$
evolve under $\step[s]{}{}$
with finality test $\isfinalf[s]$.
A \emph{backward simulation} relates a target configuration $t$
to a source configuration $s$,
ensuring that every execution step of the target
is matched by the source.
Simulations are parameterized by a postcondition $\Phi$,
which specifies the invariant expected between the final values and
memories when both programs terminate\footnote{
  This corresponds to the relation $\Phi$ in the behavioral preorder $\behord$.
}

For terminating executions,
the simulation requires that final configurations yield related values
and memories.
We formalize this condition using the predicate $\bfinalf{\Phi}$.
The assertion $\bfinal{\Phi}{t}{s}$
holds when both target and source configurations have reached final states,
producing values $v_t, v_s$ and memories $m_t, m_s$
that satisfy $\Phi(v_t, v_s, m_t, m_s)$.
This termination condition is expressed in the rule
\textsc{Lock-Final} in \cref{fig:lock-rules}.
When the source program encounters undefined behavior,
it enters a stuck state in a non-final configuration.
In this case,
\textsc{Lock-Stuck} in \cref{fig:lock-rules} immediately validates the simulation,
leaving the target program free to perform arbitrary transitions or diverge.
This asymmetry is allowed by the flexible refinement we want to prove.

For non-final configurations,
the relation must ensure that the target state is capable of progressing
and that every transition it performs corresponds to an execution step
in the source program.
Because executions may be non-terminating,
this condition is expressed coinductively in rule \textsc{Lock-Step}
in \cref{fig:lock-rules}.
Under non-determinism,
the target and source configurations are treated asymmetrically:
\begin{itemize}
  \item The target program exhibits \emph{demonic} non-determinism:
        the simulation must account for all reachable successor configurations
        $t'$.

  \item The source program exhibits \emph{angelic} non-determinism:
        the proof only needs to exhibit the existence of a single matching
        successor $s'$.
\end{itemize}
This asymmetry matches the quantification structure of the direct refinement:
we must account for every non-deterministic path chosen by the target
and justify it with an admissible path in the source.

\begin{figure}[htb]
  \begin{align*}
    &
      \begin{prooftree}
        \hypo{\bfinal{\Phi}{t}{s}}
        \infer1[\textsc{Lock-Final}]{\lsim{t}{s}{\Phi}}
      \end{prooftree}
    &
    &
      \begin{prooftree}
        \hypo{\stuck[s]{s}}
        \infer1[\textsc{Lock-Stuck}]{\lsim{t}{s}{\Phi}}
      \end{prooftree}
  \end{align*}
  \begin{equation*}
    \begin{prooftree}
      \hypo{\canp[t]{t}}
      \hypo{\forall t',
        \step[t]{t}{t'} \implies
        \exists s',
        \step[s]{s}{s'} \land \lsim{t'}{s'}{\Phi}
      }
      \infer[coind]2[\textsc{Lock-Step}]{\lsim{t}{s}{\Phi}}
    \end{prooftree}
  \end{equation*}

  \caption{Rules of the \emph{lock-step simulation}.}
  \label{fig:lock-rules}
\end{figure}

The relation $\sim_l$
defined in \cref{fig:lock-rules}
is a \emph{lock-step} backward simulation:
the target and source programs advance at the exact same pace.
While lock-step simulations suffice for simple transformations
that preserve step counts,
realistic compiler passes break this one-to-one correspondence.
An optimization pass such as branch tunneling
removes redundant jump instructions,
causing the source program to execute steps
that have no counterpart in the target.
Conversely,
lowering transformations decompose a single high-level source instruction
into multiple low-level target instructions,
requiring the target to execute several steps before the source advances.

To accommodate passes that alter step counts,
one must allow \emph{asynchronous} execution steps,
where one program takes a step while the other remains stationary.
Such a transition is called a \emph{stuttering} step.
A naive coinductive formulation might introduce unrestricted stuttering rules:
\begin{equation*}
  \begin{prooftree}
    \hypo{\canp[t]{t}}
    \hypo{\forall t',
      \step[t]{t}{t'} \implies t' \sim s
    }
    \infer[coind]2{t \sim s}
  \end{prooftree}
  \qquad\qquad
  \begin{prooftree}
    \hypo{\step[s]{s}{s'}}
    \hypo{t \sim s'}
    \infer[coind]2{t \sim s}
  \end{prooftree}
\end{equation*}

However,
adding these unconstrained rules invalidates the soundness
of the simulation relation:
one cannot prove that $t \sim s \implies t \sqsubseteq s$.
Without restrictions,
a target program could perform infinitely many stuttering steps
without the source program ever progressing.
Under such a definition,
an always diverging target program could be related to any source state.
Conversely,
a source program could perform infinitely many stuttering steps
without the target program ever progressing.
This would similarly allow relating any target program
to a diverging source state.

To ensure soundness,
the stuttering must remain finite.
This is achieved by equipping the simulation
with a well-founded measure $(W, \prec)$.
The relation is indexed by a measure $i \in W$,
written $\esim{t}{i}{s}{\Phi}$.
Whenever a stuttering step occurs,
the index strictly decreases according to $\prec$.
When both programs step simultaneously,
progress is established and the measure may be freely reset.

Because stuttering is bounded by an explicit measure,
this relation is called the \emph{explicit simulation}
\parencite{cho-23-stuttering}.
The inference rules defining it are given in
\cref{fig:explicit-sim}.
To simplify the notation for stuttering transitions,
we write $\esim{t}{\prec i}{s}{\Phi}$ as syntactic sugar for:
\begin{equation*}
  \esim{t}{\prec i}{s}{\Phi} \triangleq \exists i', i' \prec i \land \esim{t}{i'}{s}{\Phi}
\end{equation*}

The rules \textsc{ESourceSteps} and \textsc{ETargetSteps}
govern asynchronous execution steps:
\begin{itemize}
  \item In \textsc{ESourceSteps},
        the source program takes a step from $s$ to $s'$
        while the target program stutters.
        Due to angelic non-determinism in the source,
        it suffices to exhibit a single matching successor $s'$
        with a strictly decreased measure $\esim{t}{\prec i}{s'}{\Phi}$.

  \item In \textsc{ETargetSteps},
        the target $t$ must be able to make progress,
        and because target non-determinism is demonic,
        the simulation at a decreased measure $\esim{t'}{\prec i}{s}{\Phi}$
        must be established for every transition $\step[t]{t}{t'}$.
\end{itemize}
Since $(W, \prec)$ is well-founded,
an infinite sequence of stuttering steps is impossible,
guaranteeing that explicit simulation implies behavioral refinement.
However,
\textsc{EBothSteps} is the only rule
that permits resetting the stuttering measure.
This requires synchronizing the target and source steps that
justify divergence,
imposing a proof burden that we address in the following chapter.

\begin{figure}[htb]
  \begin{align*}
    &
      \begin{prooftree}
        \hypo{\bfinal{\Phi}{t}{s}}
        \infer1[\textsc{ERelated}]{\esim{t}{i}{s}{\Phi}}
      \end{prooftree}
    &
    &
      \begin{prooftree}
        \hypo{\stuck[s]{s}}
        \infer1[\textsc{EStuck}]{\esim{t}{i}{s}{\Phi}}
      \end{prooftree}
    &
    &
      \begin{prooftree}
      \hypo{\step[s]{s}{s'}}
      \hypo{\esim{t}{\prec i}{s'}{\Phi}}
      \infer[coind]2[\textsc{ESourceSteps}]{\esim{t}{i}{s}{\Phi}}
    \end{prooftree}
  \end{align*}
  \begin{gather*}
    \begin{prooftree}
      \hypo{\canp[t]{t}}
      \hypo{\forall t',
        \step[t]{t}{t'} \implies
        \esim{t'}{\prec i}{s}{\Phi}
      }
      \infer[coind]2[\textsc{ETargetSteps}]{\esim{t}{i}{s}{\Phi}}
    \end{prooftree}
    \\[1em]
    \begin{prooftree}
      \hypo{\canp[t]{t}}
      \hypo{\forall t',
        \step[t]{t}{t'} \implies
        \exists s', \exists i',
        \step[s]{s}{s'} \land \esim{t'}{i'}{s'}{\Phi}
      }
      \infer[coind]2[\textsc{EBothSteps}]{\esim{t}{i}{s}{\Phi}}
    \end{prooftree}
  \end{gather*}
  \caption{Derivation rules of the \emph{explicit simulation}.}
  \label{fig:explicit-sim}
\end{figure}

\end{document}

% LocalWords:  Coinductive coinductive asymmetricity arities
